\section{Benchmarks: numerical Chan and two references}
\label{sec:experiments}

The comparison uses three implementations throughout:
\begin{itemize}
\item \textbf{Numerical Chan (Numba):} the implementation described in this
report, with the local compression policy covered by Theorem~\ref{thm:main}.
\item \textbf{Chan $d/3$ (Python):} the repository's implementation following
Chan's functional formulation, with the same theoretical exponent.
\item \textbf{Chan $d/2$ (Python):} the simpler geometric algorithm, included
as a practical baseline despite its larger exponent.
\end{itemize}
These are comparisons of complete implementations. The numerical kernels
are compiled, while the two references run in Python; a timing advantage
therefore does not establish a better asymptotic algorithm.
Older implementation variants are documented in the repository archive
linked in Appendix~\ref{sec:archived}.

\subsection{Inputs and timing protocol}
Here $n$ is the number of input points, each defining one anchored box.
We use two families of positive, mutually nondominated points. Sphere
coordinates are absolute Gaussian samples plus $0.01$, normalized to
Euclidean length one. Simplex coordinates are exponential samples plus
$0.01$, normalized to sum to one. There is one seeded sample per family
and dimension; smaller inputs are prefixes of the largest sample.
The sphere family uses $n=64,128,256,512,1024$ in 4D--6D, plus $4096$ in
4D. The simplex family uses $n=64,256,1024$ in each dimension.

Each fresh worker runs one first call followed by three timed calls,
with a new solver on each call. The tables report the median of the
three timed calls in \emph{seconds}. Input conversion and preprocessing
are included; imports and the first call are excluded. Existing
compilation caches are retained. Compression uses the default settings;
there is no dominance prefilter, fast-math, or concurrent numerical worker.

The 45-second limit applies to the whole worker, including imports and
all four calls. Thus TO means the worker exceeded that budget before completing all
three timed calls; it is not a 45-second lower bound on one call. NR means not run after two smaller-size timeouts in the same
solver/family/dimension series. Partial samples remain available in the
raw data but are not presented as complete medians.

The environment is Windows 11 on an AMD Ryzen 7 5800H with 16 logical
CPUs, Python 3.12.14, NumPy 2.5.3 and Numba 0.68.0. CPU affinity was not
fixed and ordinary system load was uncontrolled. The measured solver
sources are unchanged from revision \texttt{095fa73}; the archived hashes
identify the exact files and inputs.

\subsection{Results and their scope}
Tables~\ref{tab:larger-sphere} and~\ref{tab:larger-simplex} retain every
planned input size. Across the three displayed methods there are 75
algorithm/input rows: 56 complete workers, 18 timeouts, and one not run.
Numerical Chan completes 18 of its 25 rows.

\begin{table}[p]
\centering\small
\setlength{\tabcolsep}{5pt}
\begin{tabular}{rrrrr}
\toprule
$d$ & $n$ & \shortstack{Numerical Chan\\Numba} & \shortstack{Chan $d/3$\\Python} & \shortstack{Chan $d/2$\\Python} \\
\midrule
4 & 64 & 0.0518 & 0.1243 & 0.0245 \\
4 & 128 & 0.1431 & 0.3472 & 0.0684 \\
4 & 256 & 0.3427 & 0.8913 & 0.1622 \\
4 & 512 & 0.8485 & 2.1206 & 0.3859 \\
4 & 1024 & 3.4890 & 4.7949 & 0.8810 \\
4 & 4096 & TO & TO & 4.5595 \\
5 & 64 & 0.2208 & 0.6844 & 0.0878 \\
5 & 128 & 0.6220 & 1.9354 & 0.2568 \\
5 & 256 & 1.9115 & 6.0686 & 0.8202 \\
5 & 512 & 5.2159 & TO & 2.0171 \\
5 & 1024 & TO & TO & 5.0315 \\
6 & 64 & 0.7754 & 3.0187 & 0.2772 \\
6 & 128 & 3.0644 & 10.9978 & 1.2158 \\
6 & 256 & 10.3680 & TO & 3.9232 \\
6 & 512 & TO & TO & TO \\
6 & 1024 & TO & NR & TO \\
\bottomrule
\end{tabular}
\caption{Sphere fronts at larger input sizes: warm median \emph{seconds}. All three methods use the same inputs. TO and NR are defined in the text.}
\label{tab:larger-sphere}
\end{table}

\begin{table}[p]
\centering\small
\setlength{\tabcolsep}{5pt}
\begin{tabular}{rrrrr}
\toprule
$d$ & $n$ & \shortstack{Numerical Chan\\Numba} & \shortstack{Chan $d/3$\\Python} & \shortstack{Chan $d/2$\\Python} \\
\midrule
4 & 64 & 0.0589 & 0.1566 & 0.0262 \\
4 & 256 & 0.3693 & 0.9423 & 0.1621 \\
4 & 1024 & 4.0905 & 5.4955 & 0.9622 \\
5 & 64 & 0.2168 & 0.7958 & 0.1137 \\
5 & 256 & 2.1954 & 7.7507 & 0.9182 \\
5 & 1024 & TO & TO & 6.1703 \\
6 & 64 & 0.8562 & 3.2904 & 0.3552 \\
6 & 256 & TO & TO & 5.2442 \\
6 & 1024 & TO & TO & TO \\
\bottomrule
\end{tabular}
\caption{Simplex fronts at larger input sizes: warm median \emph{seconds}. All three methods use the same inputs. TO and NR are defined in the text.}
\label{tab:larger-simplex}
\end{table}

On the 16 pairs with complete medians for both Numerical Chan and Chan
$d/3$, Numerical Chan is faster in all 16; the median ratio of reference
time to numerical time is $2.88$. Chan $d/2$ is faster on all 18 complete
pairs with Numerical Chan. These summaries omit incomplete runs and
must not be extrapolated to larger inputs.

The 4D sphere with $n=4096$ illustrates this limitation: Numerical Chan
returns its first value in $42.565$ seconds, compared with $23.367$
seconds for Chan $d/3$. Neither completes the three timed repetitions.
These are first-call observations, possibly including native-kernel
initialization and cache loading, not warm medians or empty-cache
compilation measurements.

The displayed methods return 242 first-call and subsequent values.
All agree with an available Chan $d/2$ value within the comparison
tolerance; the maximum scaled error is $3.3307\times10^{-16}$, using
$|v-v_{\mathrm{ref}}|/\max(1,|v_{\mathrm{ref}}|)$. No returned value is
unverified. These larger cases use cross-implementation checks, not
an exponential exact oracle. The separate small-instance checks are
described in Section~\ref{sec:software}.

Four completed Numerical Chan inputs execute compression. Among its
complete workers, the slowest/fastest warm-call ratio has median $1.05$
and maximum $1.26$. With only three repetitions and one sample per
family/dimension, the experiment does not establish an asymptotic
exponent or universal practical superiority. The complexity statement
rests on Theorem~\ref{thm:main}.

Full coordinates, individual timings, partial calls, compression counters,
and source hashes remain in \texttt{benchmarks/larger\_4d\_6d/results}.
The \href{https://github.com/emmerichmtm/FastHVChan/blob/main/NumericalChanHVND/benchmarks/larger_4d_6d/RESULTS.md}{benchmark summary}
uses the same three-method comparison as this report.
\clearpage
